\begin{abstract}
\setlength{\parindent}{2em}
\indent As lunar colonization programs transition into the phase of substantive construction, the development of an efficient, economical, and environmentally sustainable Earth-Moon logistics framework has emerged as a critical challenge. Aiming at the objective of establishing a lunar colony for 100,000 residents by 2050, this paper constructs a series of multi-dimensional coupled mathematical models to address the transport of 100 million tons of infrastructure materials, adaptation to non-ideal operating conditions, annual water resource replenishment, and Earth’s atmospheric environmental constraints.

\indent Regarding \textbf{Problem I}, a Hybrid Capacity Trade-off Optimization model based on Pareto optimality (HCTO-PM) was developed. By incorporating Wright's Law, the model establishes a dynamic cost iteration mechanism to quantitatively compare the space elevator (Option A), traditional rockets (Option B), and hybrid transportation (Option C). The results indicate that the hybrid approach performs optimally on the "time-cost" Pareto frontier: under a balanced decision-making strategy, completing the infrastructure mission requires 77 years at a total cost of approximately \$11.5 trillion.

\indent Regarding \textbf{Problem II}, a fault quantification and correction model was established. Considering stochastic factors such as cable oscillations, equipment failure rates, and rocket launch failure rates (5\%), the Pareto frontier was recalculated. Sensitivity analysis based on the single-parameter local perturbation method demonstrates that the pure elevator scheme exhibits the highest sensitivity to schedule fluctuations (with a 30.7\% duration extension), whereas the hybrid scheme demonstrates superior engineering resilience and risk mitigation through the synergistic effect of "fixed orbits plus discrete backups."

\indent Regarding \textbf{Problem III}, a Dynamic Mass Balance (DMS) model was constructed. By integrating data from human metabolism, Controlled Ecological Life Support Systems (BLSS), industrial oxygen production, and technical leakage, the net annual water resource deficit for a 100,000-person colony was calculated to be 35.4 kilotonnes. A comparison of replenishment pathways identifies the "elevator-primary with rocket-elastic-backup" mode as the optimal strategy, which reduces annual operating expenditures by approximately 14.2\% compared to pure rocket replenishment.

\indent Regarding \textbf{Problem IV}, this paper defines an Environmental Impact Unit (EIU) based on the Life Cycle Assessment (LCA) methodology. Utilizing the theory of stratospheric aerosol steady-state loading, it physically derives a threshold of 4,000 annual rocket launches as the maximum safety limit for Earth's atmosphere. By embedding this hard constraint into the optimization model, a three-dimensional Pareto mapping of "Time-Cost-Environment" was generated. The research proves that maximizing the utilization of space elevators is the only viable path for Earth's ecosystem to sustain the transport of 100 million tons of materials between the Earth and the Moon.

\indent Finally, a comprehensive policy recommendation letter was submitted to the MCM organization, covering site selection strategies, In-Situ Resource Utilization (ISRU), and dynamic scheduling mechanisms.
\begin{keywords}
Pareto multi-objective optimization; Wright's Law; Life Cycle Assessment; Space Elevator
\end{keywords}
\end{abstract}